\section{Technical Overview}
\label{sec:technical-overview}

This section gives the proof ideas in the same order as the main text. Full formal versions appear in Appendices~\ref{app:refinement-proofs}, \ref{app:value-proofs}, and \ref{app:algorithm-bridge-proofs}.

\begin{theorem}[Minimal finite refinement, proof sketch for Theorem~\ref{thm:minimal-refinement:formal}]
\label{thm:minimal-refinement:sketch}
The informal statement is Theorem~\ref{thm:minimal-refinement:informal}.
\end{theorem}
\begin{proof}[Proof sketch]
Every atom of $\calP_E$ lies inside an atom of $\calP$, so $\calP_E$ refines $\calP$. Also
\begin{align*}
E=\bigcup_{B\in\calP}(B\cap E),
\end{align*}
so $E$ is $\calP_E$-measurable. If $\calQ$ refines $\calP$ and makes $E$ measurable, then every atom $Q\in\calQ$ satisfies
\begin{align*}
Q\subseteq B
\end{align*}
for some $B\in\calP$, and also either $Q\subseteq E$ or $Q\subseteq E^c$. Hence
\begin{align*}
Q\subseteq B\cap E\qquad \text{or}\qquad Q\subseteq B\cap E^c.
\end{align*}
Thus $\calQ$ refines $\calP_E$.
\end{proof}

\begin{theorem}[Minimal sigma-algebra refinement, proof sketch for Theorem~\ref{thm:sigma-refinement:formal}]
\label{thm:sigma-refinement:sketch}
The informal statement is Theorem~\ref{thm:sigma-refinement:informal}.
\end{theorem}
\begin{proof}[Proof sketch]
By definition,
\begin{align*}
\calF\vee\sigma(Z)=\sigma(\calF\cup\sigma(Z)).
\end{align*}
This sigma-algebra contains $\calF$ and makes $Z$ measurable. If $\calG$ also contains $\calF$ and makes $Z$ measurable, then
\begin{align*}
\calF\cup\sigma(Z)\subseteq\calG.
\end{align*}
Since $\calG$ is a sigma-algebra,
\begin{align*}
\sigma(\calF\cup\sigma(Z))\subseteq\calG.
\end{align*}
\end{proof}

\begin{theorem}[Contradiction, proof sketch for Theorem~\ref{thm:contradiction:formal}]
\label{thm:contradiction:sketch}
The informal statement is Theorem~\ref{thm:contradiction:informal}.
\end{theorem}
\begin{proof}[Proof sketch]
If $f$ is $\calP$-measurable and $\omega,\omega'$ lie in the same atom, then
\begin{align*}
f(\omega)=f(\omega').
\end{align*}
If $f(\omega)=y(\omega)$ and $f(\omega')=y(\omega')$, then
\begin{align*}
y(\omega)=y(\omega'),
\end{align*}
contradicting the hypothesis.
\end{proof}

\begin{theorem}[Risk monotonicity, proof sketch for Theorem~\ref{thm:risk-monotonicity:formal}]
\label{thm:risk-monotonicity:sketch}
The informal statement is Theorem~\ref{thm:risk-monotonicity:informal}.
\end{theorem}
\begin{proof}[Proof sketch]
If $\calF\subseteq\calG$, then every $\calF$-measurable action is $\calG$-measurable. Therefore
\begin{align*}
\{a:a\text{ is }\calF\text{-measurable}\}\subseteq\{a:a\text{ is }\calG\text{-measurable}\}.
\end{align*}
Taking the infimum of the same expected loss over the larger set gives
\begin{align*}
R_L(\calG)\le R_L(\calF).
\end{align*}
\end{proof}

\begin{theorem}[Log-loss value, proof sketch for Theorem~\ref{thm:log-value:formal}]
\label{thm:log-value:sketch}
The informal statement is Theorem~\ref{thm:log-value:informal}.
\end{theorem}
\begin{proof}[Proof sketch]
For any $\calF$-measurable conditional distribution $q$,
\begin{align*}
\E[-\log q(Y)]=H(Y\mid\calF)+\E\left[D_{\mathrm{KL}}(\mathbb P(Y\in\cdot\mid\calF)\|q(\cdot))\right].
\end{align*}
The minimizer is $q^\star=\mathbb P(Y\in\cdot\mid\calF)$, so $R_{\log}(\calF)=H(Y\mid\calF)$. Similarly $R_{\log}(\calG)=H(Y\mid\calG)$. Therefore
\begin{align*}
R_{\log}(\calF)-R_{\log}(\calG)=H(Y\mid\calF)-H(Y\mid\calG)=I(Y;\calG\mid\calF).
\end{align*}
\end{proof}

\begin{theorem}[Squared-loss value, proof sketch for Theorem~\ref{thm:squared-value:formal}]
\label{thm:squared-value:sketch}
The informal statement is Theorem~\ref{thm:squared-value:informal}.
\end{theorem}
\begin{proof}[Proof sketch]
The optimal $\calF$-predictor is $\E[Y\mid\calF]$. Since $\calF\subseteq\calG$,
\begin{align*}
Y-\E[Y\mid\calF]=\left(Y-\E[Y\mid\calG]\right)+\left(\E[Y\mid\calG]-\E[Y\mid\calF]\right).
\end{align*}
The two terms are orthogonal because the first has conditional mean zero given $\calG$ and the second is $\calG$-measurable. Taking squared $\ell_2$ norms and expectations gives the identity.
\end{proof}

\begin{theorem}[Prototype edge separation, proof sketch for Theorem~\ref{thm:prototype-edge:formal}]
\label{thm:prototype-edge:sketch}
The informal statement is Theorem~\ref{thm:prototype-edge:informal}.
\end{theorem}
\begin{proof}[Proof sketch]
With $u=(c_i-c_j)/\norm{c_i-c_j}_2$ and $b=u^\top(c_i+c_j)/2$,
\begin{align*}
u^\top c_i-b&=\frac{1}{2}u^\top(c_i-c_j)=\frac{1}{2}\norm{c_i-c_j}_2,\\
u^\top c_j-b&=-\frac{1}{2}u^\top(c_i-c_j)=-\frac{1}{2}\norm{c_i-c_j}_2.
\end{align*}
\end{proof}

\begin{theorem}[DGM data structure, proof sketch for Theorem~\ref{thm:dpm-complexity:formal}]
\label{thm:dpm-complexity:sketch}
The informal statement is Theorem~\ref{thm:dpm-complexity:informal}.
\end{theorem}
\begin{proof}[Proof sketch]
A prediction computes distances from $h$ to $N_t$ prototypes in $\R^d$, costing
\begin{align*}
O(N_td).
\end{align*}
It then evaluates $k$ selected local memories, costing $O(kd+C_{\mathrm{local}})$ including sparse routing weights. Repair adds the local repair cost. Creating one prototype stores a vector in $\R^d$. Adding $q$ distinction edges stores $q$ vectors and thresholds, costing
\begin{align*}
O(qd).
\end{align*}
For embedding-mean memories, each prototype stores $c_j\in\R^d$ and $r_j\in\R^{d_e}$, while each edge stores one distinction vector in $\R^d$. Hence memory is
\begin{align*}
O(N_t(d+d_e)+|E_t|d).
\end{align*}
\end{proof}

\begin{theorem}[DGM layer cost, proof sketch for Theorem~\ref{thm:dgm-layer-cost:formal}]
\label{thm:dgm-layer-cost:sketch}
The informal statement is Theorem~\ref{thm:dgm-layer-cost:informal}.
\end{theorem}
\begin{proof}[Proof sketch]
Computing $q=Qh$ costs $O(dd_q)$. Retrieval costs $C_{\mathrm{ret}}(N_t,d_q)$ by definition. For each selected prototype, a rank-$r$ adapter computes
\begin{align*}
B_jA_jh.
\end{align*}
The two matrix-vector products cost $O(dr)$, so $k$ selected adapters cost
\begin{align*}
O(kdr).
\end{align*}
The detached write path repeats query and retrieval, then updates only selected local states. A rank-one local repair or low-rank sufficient-statistic update costs $O(dr)$ per selected memory, with $O(d_q)$ additional prototype bookkeeping. Hence write overhead is
\begin{align*}
O(dd_q+C_{\mathrm{ret}}(N_t,d_q)+k(d_q+dr)).
\end{align*}
Because the write path is detached, backpropagation does not need intermediate copies of $\calS_t$ for all $t$. It stores the selected read computation used in the forward loss.
\end{proof}

\begin{theorem}[One-hot memory, proof sketch for Theorem~\ref{thm:one-hot-measurable:formal}]
\label{thm:one-hot-measurable:sketch}
The informal statement is Theorem~\ref{thm:one-hot-measurable:informal}.
\end{theorem}
\begin{proof}[Proof sketch]
If $\omega,\omega'\in B_i$, then both one-hot features equal $e_i$, so
\begin{align*}
W\phi_\calP(\omega)=We_i=W\phi_\calP(\omega').
\end{align*}
Thus $W\phi_\calP$ is constant on atoms. Conversely, if $f$ is constant on each $B_i$, put that constant value in column $i$ of $W$. Then
\begin{align*}
f(\omega)=W\phi_\calP(\omega)
\end{align*}
for every $\omega$.
\end{proof}

\begin{theorem}[CPM conflict, proof sketch for Theorem~\ref{thm:cpm-conflict-distinction:formal}]
\label{thm:cpm-conflict-distinction:sketch}
The informal statement is Theorem~\ref{thm:cpm-conflict-distinction:informal}.
\end{theorem}
\begin{proof}[Proof sketch]
If $\omega,\omega'\in B_i$, then
\begin{align*}
\phi_\calP(\omega)=\phi_\calP(\omega')=e_i.
\end{align*}
After committing $We_i=v(\omega)$, the later evidence $We_i=v(\omega')$ has the same key. Its innovation relative to the old key span is
\begin{align*}
z=(I-e_ie_i^\top)e_i=0.
\end{align*}
Its residual is
\begin{align*}
e=v(\omega')-We_i=v(\omega')-v(\omega)\ne0.
\end{align*}
\end{proof}
